\subsection{Borel 图像定理}

定理 1. --- 设 E 是一个作为 Banach 空间的归纳极限的局部凸空间，F 是一个 Souslin 局部凸空间，例如一个 Lusin 空间（GT, IX, § 6, No.~2 与 No.~4），u 是从 E 到 F 中的一个线性映射。若 u 的图像是 E × F 的一个 Borel 子集，则 u 连续。

设 \(E_i\) 是一族 Banach 空间，\((u_i)\) 是一族连续线性映射 \(u_i: E_i \to E\) ，使得 \(E\) 的拓扑是使诸 \(u_i\) 连续的最细局部凸拓扑。只需证明诸复合映射 \(u \circ u_i\) 都是连续的，或者事实上（GT, IX, § 2, No.~6, 命题 10）证明 \(u \circ u_i\) 到每个满足第一可数公理的闭子空间 \(G\) （在 \(E_i\) 中）上的限制都是连续的。这一限制的图像是 \(u\) 的图像在连续映射 \(u_i \times \text{Id}_F: G \times F \to E \times F\) 下的原像，因而是 \(G \times F\) 中的一个 Borel 集。此外，\(G \times F\) 是一个 Souslin 空间，而 Souslin 空间的每个 Borel 子集都是 Souslin 空间（GT, IX, § 6, No.~3, 命题 10）。于是定理 1 由 GT, IX, § 6, No.~8 的定理 4 得出。

注记. --- 回忆（III, 第 12 页），每个同调的 Hausdorff 半完备空间，例如每个 Fréchet 空间，都是 Banach 空间的归纳极限。* 这对自反 Fréchet 空间的强对偶也同样成立（IV, 第 23 页, 命题 4）。*

